% !TeX root = ../papers/nc-p3.tex
\section{三重点の四次鉛筆と復元次数二}
\label{sec:configuration-quartic-pencil}
$D=k^{20}$に成分ごとの積を入れ、$Q\subset D$を十次元とする。
座標評価を$p_i\in Q^*$とし、$H(Q)$を各$[p_i]$で三重以上に消える
$\mathbb P(Q^*)$の四次式の線型空間とする。以下、標数零を用い、
有限体での検算には標数$> 4$を用いる。

\begin{proposition}\label{prop:quartic-root-pencil}
$L\in\Gr(4,D)$が$L^2\subset Q$を満たすと、
$H(Q)$に自然な非零の射影元$[F_L]$がある。各$i$でその対称四重線型形式の切断
\[
 Q^*\longrightarrow \operatorname{Sym}^2 Q,\qquad
 v\longmapsto F_L(p_i,v,-,-)
\]
は階数6以下である。$L^2=Q$なら$F_L$の三重点軌跡は二次Veronese三次元多様体で、
$F_L$が一致する二つの全階数平方根は、対応する二十点で独立な行単元倍を除いて同じである
（各$p_i\ne0$とする）。
\end{proposition}
\begin{proof}
$\mu:\operatorname{Sym}^2L\to Q$を積として
$F_L(v)=\det(\mu^*v)$と取る。$\mu^*p_i$は$L$の第$i$行の平方である。
独立な四行の平方の和は非退化なので、$\mu$の全射性を仮定しなくても
この四次式は恒等的に零にならない。階数1の対称四行行列での最低次項は
核三平面上の対称行列の行列式である。その三次式が使う変数は六つであり、
零行に対応する階数0では三次項自身が零となる。制限写像によって変数の個数は増えない。
全階数なら三重以上の軌跡はちょうど階数1の射影行列の全体である。
この三次元多様体を保つ射影自己同型は元の$\mathbb P^3$の自己同型から誘導される。
従って二つの表示の各行は共通の基底変更の後に比例する。さらに双方の平方が$Q$で、
一方の三乗が$D$に同型なら、命題\ref{prop:finite-square-ratio}より
行単元の平方はスカラーで、違いは全体倍率と各行の符号に限られる。
\end{proof}

\begin{theorem}\label{thm:exact-two-configurations}
十五次元成分の幾何学的一般点と反対階数16の有理例$t$で、
第一配置からの添字保持同型類の復元次数はちょうど2である。
四辺が有限な標準AS延長は元の窓とその逆向きの窓で尽きる。
第一辺まで指定すると延長は一意である。反対階数10の有理例$r$の
第一配置でも同じ二つの向きで尽きる。
\end{theorem}
\begin{proof}
第一位置の行列の列空間を$A$とし、$Q=A^2$とする。
順に他の三位置を$B,C,E$と書く。辺上の点の番号を共有し、
定理\ref{thm:finite-square-chart}の完全交叉のtraceを使って
$C^2=Q$, $B^2=E^2=Q^\perp$とできる。後者では$D$の標準内積を使い、
行ごとの直線束の自明化を選ぶ。ここで扱う基点と一般点では四移動は恒等である。実際、二十点の番号付け上の戻り置換は局所一定で、両有理例で恒等と検算されている。従って$Q^\perp$の射影評価点は$B$の各行の
二次Veronese像に射影同値である。$t$の反対位置の点配置は互いに射影非同値
なので、両側で少なくとも二つの異なる行列式四次式がある。

計算を記す。基底の十変数の四次単項式は715個あり、各評価点で
$F(v+s p_i)$の$s^2$係数（$v$の二次式）を零と置く。
対称化により、これは三重以上の条件と同値である。二十点を合わせた
1100行715列の行列の核を取り、その基底$F_0,F_1$の各々から
最初の点の上記切断を作る。切断は55行10列であり、$F_0+uF_1$に対して
その七次小行列式の線型結合を二本取る。具体的には左右から七行と七列に圧縮した
行列式であり、各々次数7以下なので八値で正確に計算できる。
\path{computations/verify_configuration_quartic_pencil.py}の基底と圧縮を使った結果は次の通りである。
\[
\begin{array}{c|c|c|c|c}
 \text{基点}&\text{素数}&\text{平方の位置}&\dim H&\text{二本の最大公約式}\\\hline
 t&2657&A&2&u(u+277)\\
 t&2657&B&2&u(u+22)\\
 r&12619&A&2&u^2\\
 r&12619&B&2&u(u+2481)
\end{array}
\]
どれも無限遠で少なくとも一方が非零である。この計算が与えるのは、
可能な行列式を含む次数2の上界である。一般の鉛筆の全ての元が行列式になることは仮定しない。

$t$の近くで番号をエタール局所的に取り、二つの既知の行列式により
核次元は少なくとも2となる。上の小行列で同時に高々2なので鉛筆は局所自由である。
有限体から標数零の$t$へも同じ議論が使える。射影鉛筆上の二本の七次式の
共通零点は特殊ファイバーで長さ2なので、近傍で有限かつ長さは高々2となる。
これは射影性、有限射の直像に対する半連続性による。よって$t$と一般点で
可能な行列式は両側とも既知の二本だけである。これらは十変数全てを必要とする。
実際、行列式の三重軌跡は滑らかなVeronese像であり、非零な方向を変数から
除ければその射影点で四重になってしまう。従って$\mu$が全射でない平方根から
同じ式を得ることもない。命題\ref{prop:quartic-root-pencil}により、退化解も含め
Grassmannスキームの全幾何点が各側ちょうど二つの符号軌道で尽くされる。
既知の解は孤立被約で、一軌道は$2^{19}$点である。

こうして第二空間の候補は$B,E$、第三空間の候補は$A,C$に限られる。
$B$の後に$A$を取る候補は前節の三次交叉で排除される。$E$を先にした
逆向きの窓にも同じ検査を行うと$t$で交叉は零であり、一般点でも零である。
各組は連続二辺の積の核を決め、そのAS延長は系\ref{cor:two-window-extension}
により一意である。二つの向きは第二配置が異なるため同型でない。
特に十五次元成分の関数体は第一配置の像の関数体上二次であり、非自明な
自己同型は逆転である。ここで全四辺有限の他成分から同じ配置に来るAS窓も既に数えた。

$r$では奇側$B,E$が異なるので同じ証明が適用できる。偶側は標数12619で
上の共通零点が一つしかなく、既知の$A$の行列式がその点を占める。
$A$の三次評価は階数20であるから、上の議論により全幾何点は$A$の符号倍である。
前節の検算はこの剰余体でも各点の接次元1、二次障害非零を与えるので、
その局所長さは2で、全長さは$2^{20}$である。平方根Grassmannスキームは
固有なので有理点$r$の標数零ファイバーも有限で長さは高々$2^{20}$である。
前節で証明した$A$の長さ2の符号軌道が既にこれを尽くす。従って第三空間は
$A$だけとなる。二つの第二空間は実際に正則な両方向を与えるため結論を得る。
\end{proof}
